\documentclass[11pt]{article}

\usepackage{amsmath, amssymb, amsthm}
\usepackage{mathtools}
\usepackage{geometry}
\usepackage{enumitem}
\usepackage{hyperref}

\geometry{margin=1in}

\newtheorem{theorem}{Theorem}[section]
\newtheorem{proposition}[theorem]{Proposition}
\newtheorem{lemma}[theorem]{Lemma}
\newtheorem{corollary}[theorem]{Corollary}
\theoremstyle{definition}
\newtheorem{definition}[theorem]{Definition}
\newtheorem{remark}[theorem]{Remark}

\DeclareMathOperator{\Cov}{Cov}
\DeclareMathOperator{\Var}{Var}
\DeclareMathOperator{\Range}{Range}
\DeclareMathOperator{\Tr}{Tr}

\newcommand{\HK}{\mathcal{H}_K}
\newcommand{\X}{\mathcal{X}}
\newcommand{\B}{\mathcal{B}}
\newcommand{\E}{\mathbb{E}}
\newcommand{\R}{\mathbb{R}}
\newcommand{\N}{\mathbb{N}}
\newcommand{\iid}{\stackrel{\mathrm{iid}}{\sim}}
\newcommand{\inprod}[2]{\langle #1, #2 \rangle}
\newcommand{\norm}[1]{\| #1 \|}

\title{Gaussian Measures from Positive Semi-Definite Kernels:\\
Formal Definitions and the RKHS Prior Correspondence}
\author{}
\date{}

\begin{document}

\maketitle

\begin{abstract}
We provide formal definitions for constructing Gaussian measures from positive semi-definite kernels and establish the precise correspondence between these measures and the informal prior ``$p(f) \propto \exp(-\frac{1}{2}\|f\|_{\HK}^2)$'' via the Onsager-Machlup functional.
\end{abstract}

\tableofcontents

%%%%%%%%%%%%%%%%%%%%%%%%%%%%%%%%%%%%%%%%%%%%%%%%%%%%%%%%%%%%%%%%%%%%%%%%%%%%%%%
\section{Setup and Notation}
%%%%%%%%%%%%%%%%%%%%%%%%%%%%%%%%%%%%%%%%%%%%%%%%%%%%%%%%%%%%%%%%%%%%%%%%%%%%%%%

Let $(\X, \mathcal{A}, P_X)$ be a probability space where $\X$ is the input domain. Let $K: \X \times \X \to \R$ be a symmetric positive semi-definite (PSD) kernel satisfying the integrability condition:
\begin{equation}
    \int_{\X} K(x, x) \, dP_X(x) < \infty.
\end{equation}

We work primarily in the Hilbert space $L^2(P_X)$ with inner product
\begin{equation}
    \inprod{f}{g}_{L^2} = \int_{\X} f(x) g(x) \, dP_X(x).
\end{equation}

%%%%%%%%%%%%%%%%%%%%%%%%%%%%%%%%%%%%%%%%%%%%%%%%%%%%%%%%%%%%%%%%%%%%%%%%%%%%%%%
\section{The Integral Operator}
%%%%%%%%%%%%%%%%%%%%%%%%%%%%%%%%%%%%%%%%%%%%%%%%%%%%%%%%%%%%%%%%%%%%%%%%%%%%%%%

\begin{definition}[Integral Operator]
The integral operator $T_K: L^2(P_X) \to L^2(P_X)$ associated with kernel $K$ is defined by
\begin{equation}
    (T_K g)(x) = \int_{\X} K(x, y) \, g(y) \, dP_X(y).
\end{equation}
\end{definition}

\begin{theorem}[Spectral Decomposition]
\label{thm:spectral}
The operator $T_K$ is self-adjoint, positive, and trace-class. It admits a spectral decomposition
\begin{equation}
    T_K \phi_k = \mu_k \phi_k, \quad k \in \N,
\end{equation}
where:
\begin{enumerate}[label=(\roman*)]
    \item $\mu_1 \geq \mu_2 \geq \cdots > 0$ are the positive eigenvalues (counted with multiplicity),
    \item $\{\phi_k\}_{k=1}^\infty$ is an orthonormal basis of $\overline{\Range(T_K)} \subseteq L^2(P_X)$,
    \item $\sum_{k=1}^\infty \mu_k = \Tr(T_K) = \int_{\X} K(x,x) \, dP_X(x) < \infty$.
\end{enumerate}
The kernel admits the Mercer representation
\begin{equation}
    K(x, y) = \sum_{k=1}^\infty \mu_k \, \phi_k(x) \, \phi_k(y),
\end{equation}
with convergence in $L^2(P_X \otimes P_X)$.
\end{theorem}

%%%%%%%%%%%%%%%%%%%%%%%%%%%%%%%%%%%%%%%%%%%%%%%%%%%%%%%%%%%%%%%%%%%%%%%%%%%%%%%
\section{The Reproducing Kernel Hilbert Space}
%%%%%%%%%%%%%%%%%%%%%%%%%%%%%%%%%%%%%%%%%%%%%%%%%%%%%%%%%%%%%%%%%%%%%%%%%%%%%%%

\begin{definition}[RKHS via Spectral Representation]
\label{def:rkhs}
The reproducing kernel Hilbert space (RKHS) $\HK$ associated with $K$ is
\begin{equation}
    \HK = \left\{ f = \sum_{k=1}^\infty \theta_k \phi_k : \norm{f}_{\HK}^2 := \sum_{k=1}^\infty \frac{\theta_k^2}{\mu_k} < \infty \right\},
\end{equation}
equipped with the inner product
\begin{equation}
    \inprod{f}{g}_{\HK} = \sum_{k=1}^\infty \frac{\theta_k \eta_k}{\mu_k}, \quad \text{where } f = \sum_k \theta_k \phi_k, \; g = \sum_k \eta_k \phi_k.
\end{equation}
\end{definition}

\begin{proposition}[Operator Characterization]
\label{prop:rkhs-operator}
The RKHS admits the equivalent characterization
\begin{equation}
    \HK = \Range(T_K^{1/2})
\end{equation}
with norm
\begin{equation}
    \norm{f}_{\HK} = \norm{T_K^{-1/2} f}_{L^2}.
\end{equation}
The inclusion $\HK \hookrightarrow L^2(P_X)$ is continuous and dense in $\overline{\Range(T_K)}$.
\end{proposition}

\begin{remark}
The space $\HK$ is ``smaller'' than $L^2(P_X)$: for $f \in \HK$, the coefficients $\theta_k = \inprod{f}{\phi_k}_{L^2}$ must decay faster than $\sqrt{\mu_k}$.
\end{remark}

%%%%%%%%%%%%%%%%%%%%%%%%%%%%%%%%%%%%%%%%%%%%%%%%%%%%%%%%%%%%%%%%%%%%%%%%%%%%%%%
\section{The Gaussian Measure}
%%%%%%%%%%%%%%%%%%%%%%%%%%%%%%%%%%%%%%%%%%%%%%%%%%%%%%%%%%%%%%%%%%%%%%%%%%%%%%%

\begin{definition}[Gaussian Measure via Characteristic Functional]
\label{def:gaussian-measure}
The \emph{Gaussian measure} $\mu_K$ on $L^2(P_X)$ is the unique Borel probability measure with characteristic functional
\begin{equation}
    \boxed{
    \hat{\mu}_K(g) := \int_{L^2(P_X)} e^{i\inprod{g}{f}_{L^2}} \, d\mu_K(f) = \exp\left( -\frac{1}{2} \inprod{g}{T_K g}_{L^2} \right)
    }
\end{equation}
for all $g \in L^2(P_X)$.
\end{definition}

\begin{theorem}[Existence and Characterization]
\label{thm:gaussian-existence}
The measure $\mu_K$ exists, is unique, and satisfies the following:

\begin{enumerate}[label=(\roman*)]
    \item \textbf{Karhunen-Lo\`eve representation:}
    \begin{equation}
        \boxed{
        f \sim \mu_K \iff f = \sum_{k=1}^\infty \sqrt{\mu_k} \, \xi_k \, \phi_k, \quad \xi_k \iid N(0,1)
        }
    \end{equation}
    where the series converges in $L^2(\Omega; L^2(P_X))$.
    
    \item \textbf{Mean and covariance:}
    \begin{align}
        \E_{\mu_K}[f] &= 0, \\
        \E_{\mu_K}[\inprod{f}{g}_{L^2} \inprod{f}{h}_{L^2}] &= \inprod{g}{T_K h}_{L^2}, \quad \forall g, h \in L^2(P_X).
    \end{align}
    Equivalently, $\E_{\mu_K}[f \otimes f] = T_K$ as operators.
    
    \item \textbf{Finite-dimensional distributions:} For any $x_1, \ldots, x_n \in \X$,
    \begin{equation}
        (f(x_1), \ldots, f(x_n))^\top \sim N(0, K_{\mathbf{X}\mathbf{X}}),
    \end{equation}
    where $(K_{\mathbf{X}\mathbf{X}})_{ij} = K(x_i, x_j)$.
\end{enumerate}
\end{theorem}

\begin{remark}[Gaussian Process Equivalence]
The measure $\mu_K$ is the law of the Gaussian process $f \sim \mathcal{GP}(0, K)$.
\end{remark}

%%%%%%%%%%%%%%%%%%%%%%%%%%%%%%%%%%%%%%%%%%%%%%%%%%%%%%%%%%%%%%%%%%%%%%%%%%%%%%%
\section{The Cameron-Martin Theorem}
%%%%%%%%%%%%%%%%%%%%%%%%%%%%%%%%%%%%%%%%%%%%%%%%%%%%%%%%%%%%%%%%%%%%%%%%%%%%%%%

The following theorem establishes $\HK$ as the Cameron-Martin space of $\mu_K$, providing the fundamental link between the RKHS and the Gaussian measure.

\begin{theorem}[Cameron-Martin]
\label{thm:cameron-martin}
Let $\mu_K$ be the Gaussian measure on $L^2(P_X)$ and $\HK$ its associated RKHS. Then:

\begin{enumerate}[label=(\roman*)]
    \item \textbf{Null RKHS:} Sample paths lie outside the RKHS almost surely:
    \begin{equation}
        \mu_K(\HK) = 0.
    \end{equation}
    
    \item \textbf{Admissible shifts:} For $h \in L^2(P_X)$, define the translated measure $\mu_K^h(A) := \mu_K(A - h)$. Then
    \begin{equation}
        \mu_K^h \ll \mu_K \iff h \in \HK.
    \end{equation}
    That is, $\HK$ is exactly the set of admissible shifts preserving absolute continuity.
    
    \item \textbf{Radon-Nikodym derivative:} For $h \in \HK$,
    \begin{equation}
        \boxed{
        \frac{d\mu_K^h}{d\mu_K}(f) = \exp\left( \inprod{h}{f}_{\HK} - \frac{1}{2}\norm{h}_{\HK}^2 \right)
        }
    \end{equation}
    where the pairing $\inprod{h}{f}_{\HK}$ extends to $f \in L^2(P_X)$ via
    \begin{equation}
        \inprod{h}{f}_{\HK} = \sum_{k=1}^\infty \frac{\eta_k \theta_k}{\mu_k}, \quad h = \sum_k \eta_k \phi_k, \; f = \sum_k \theta_k \phi_k.
    \end{equation}
\end{enumerate}
\end{theorem}

\begin{corollary}
The RKHS norm $\norm{\cdot}_{\HK}$ measures the ``cost'' of deterministic shifts under the Gaussian measure.
\end{corollary}

%%%%%%%%%%%%%%%%%%%%%%%%%%%%%%%%%%%%%%%%%%%%%%%%%%%%%%%%%%%%%%%%%%%%%%%%%%%%%%%
\section{The Onsager-Machlup Functional}
%%%%%%%%%%%%%%%%%%%%%%%%%%%%%%%%%%%%%%%%%%%%%%%%%%%%%%%%%%%%%%%%%%%%%%%%%%%%%%%

The formal expression ``$p(f) \propto \exp(-\frac{1}{2}\norm{f}_{\HK}^2)$'' is made rigorous by the following theorem.

\begin{theorem}[Onsager-Machlup]
\label{thm:onsager-machlup}
Let $\mu_K$ be a Gaussian measure on a Banach space $\B \supseteq \HK$ with Cameron-Martin space $\HK$. Let $\norm{\cdot}_{\B}$ be a measurable norm satisfying appropriate regularity conditions.\footnote{Specifically, $\norm{\cdot}_{\B}$ should be measurable and the small ball probabilities $\mu_K(B_\epsilon(0))$ should satisfy certain asymptotic conditions; see Zeitouni (1989).} Then for any $f, g \in \HK$:
\begin{equation}
    \boxed{
    \lim_{\epsilon \to 0} \frac{\mu_K\bigl(\{h \in \B : \norm{h - f}_{\B} < \epsilon\}\bigr)}{\mu_K\bigl(\{h \in \B : \norm{h - g}_{\B} < \epsilon\}\bigr)} = \exp\left( -\frac{1}{2}\norm{f}_{\HK}^2 + \frac{1}{2}\norm{g}_{\HK}^2 \right)
    }
\end{equation}
\end{theorem}

\begin{corollary}[RKHS Norm as Log-Density]
Taking $g = 0$:
\begin{equation}
    \lim_{\epsilon \to 0} \frac{\mu_K(B_\epsilon(f))}{\mu_K(B_\epsilon(0))} = \exp\left( -\frac{1}{2}\norm{f}_{\HK}^2 \right).
\end{equation}
This is the precise sense in which
\begin{equation}
    ``p(f) \propto \exp\left(-\frac{1}{2}\norm{f}_{\HK}^2\right)."
\end{equation}
\end{corollary}

\begin{remark}
The function $f \mapsto \frac{1}{2}\norm{f}_{\HK}^2$ is called the \emph{Onsager-Machlup functional} of $\mu_K$. It plays the role of a ``negative log-density'' for the Gaussian measure.
\end{remark}

%%%%%%%%%%%%%%%%%%%%%%%%%%%%%%%%%%%%%%%%%%%%%%%%%%%%%%%%%%%%%%%%%%%%%%%%%%%%%%%
\section{Connection to Bayesian Inference and Kernel Ridge Regression}
%%%%%%%%%%%%%%%%%%%%%%%%%%%%%%%%%%%%%%%%%%%%%%%%%%%%%%%%%%%%%%%%%%%%%%%%%%%%%%%

\begin{theorem}[Posterior and MAP Estimator]
\label{thm:posterior}
Suppose we observe
\begin{equation}
    Y_i = f(X_i) + \varepsilon_i, \quad i = 1, \ldots, n,
\end{equation}
where $\varepsilon_i \iid N(0, \sigma^2)$, $X_i \iid P_X$, and the prior is $f \sim \mu_K$. Then:

\begin{enumerate}[label=(\roman*)]
    \item The posterior $\mu_K^{(n)}(\cdot \mid \mathbf{Y}, \mathbf{X})$ is a Gaussian measure.
    
    \item The posterior mean coincides with the kernel ridge regression estimator:
    \begin{equation}
        \boxed{
        \hat{f} = \arg\min_{f \in \HK} \left\{ \frac{1}{\sigma^2} \sum_{i=1}^n (Y_i - f(X_i))^2 + \norm{f}_{\HK}^2 \right\}
        }
    \end{equation}
    
    \item In closed form:
    \begin{equation}
        \hat{f}(\cdot) = \mathbf{k}(\cdot)^\top (K_{\mathbf{X}\mathbf{X}} + \sigma^2 I)^{-1} \mathbf{Y},
    \end{equation}
    where $\mathbf{k}(x) = (K(x, X_1), \ldots, K(x, X_n))^\top$.
\end{enumerate}
\end{theorem}

\begin{remark}
The RKHS norm appears as the regularizer in the MAP estimator precisely because it is the Onsager-Machlup functional of the GP prior. This establishes the equivalence:
\begin{center}
\emph{GP regression with covariance $K$} $\longleftrightarrow$ \emph{Kernel ridge regression with RKHS $\HK$}.
\end{center}
\end{remark}

%%%%%%%%%%%%%%%%%%%%%%%%%%%%%%%%%%%%%%%%%%%%%%%%%%%%%%%%%%%%%%%%%%%%%%%%%%%%%%%
\section{Sum Kernels}
%%%%%%%%%%%%%%%%%%%%%%%%%%%%%%%%%%%%%%%%%%%%%%%%%%%%%%%%%%%%%%%%%%%%%%%%%%%%%%%

For $K = \sum_{j=1}^J a_j K_{A_j}$ with $a_j > 0$ and each $K_{A_j}$ PSD:

\begin{theorem}[Convolution Structure]
\label{thm:sum-kernel}
The Gaussian measure satisfies
\begin{equation}
    \mu_K = \mu_{a_1 K_{A_1}} * \mu_{a_2 K_{A_2}} * \cdots * \mu_{a_J K_{A_J}},
\end{equation}
where $*$ denotes convolution of measures. Equivalently,
\begin{equation}
    \boxed{
    f \sim \mu_K \iff f = \sum_{j=1}^J f_j, \quad f_j \stackrel{\mathrm{ind}}{\sim} \mu_{a_j K_{A_j}}
    }
\end{equation}
\end{theorem}

\begin{theorem}[RKHS of Sum Kernel]
\label{thm:sum-rkhs}
The RKHS of $K = \sum_{j=1}^J a_j K_{A_j}$ is
\begin{equation}
    \HK = \mathcal{H}_{K_{A_1}} + \cdots + \mathcal{H}_{K_{A_J}} = \left\{ \sum_{j=1}^J f_j : f_j \in \mathcal{H}_{K_{A_j}} \right\}
\end{equation}
with the infimal convolution norm:
\begin{equation}
    \boxed{
    \norm{f}_{\HK}^2 = \inf_{\substack{f = \sum_{j=1}^J f_j \\ f_j \in \mathcal{H}_{K_{A_j}}}} \sum_{j=1}^J \frac{1}{a_j} \norm{f_j}_{\mathcal{H}_{K_{A_j}}}^2
    }
\end{equation}
\end{theorem}

%%%%%%%%%%%%%%%%%%%%%%%%%%%%%%%%%%%%%%%%%%%%%%%%%%%%%%%%%%%%%%%%%%%%%%%%%%%%%%%
\section{Summary: The Complete Correspondence}
%%%%%%%%%%%%%%%%%%%%%%%%%%%%%%%%%%%%%%%%%%%%%%%%%%%%%%%%%%%%%%%%%%%%%%%%%%%%%%%

\begin{center}
\fbox{\parbox{0.9\textwidth}{
\textbf{Given:} PSD kernel $K: \X \times \X \to \R$

\vspace{0.5em}
$\Downarrow$
\vspace{0.5em}

\textbf{Integral operator:} $T_K: L^2(P_X) \to L^2(P_X)$, \quad $T_K \phi_k = \mu_k \phi_k$

\vspace{0.5em}
$\Downarrow$
\vspace{0.5em}

\textbf{RKHS:} $\displaystyle \HK = \left\{ f = \sum_k \theta_k \phi_k : \sum_k \frac{\theta_k^2}{\mu_k} < \infty \right\}$

\vspace{0.5em}
$\Downarrow$
\vspace{0.5em}

\textbf{Gaussian measure:} $\displaystyle \hat{\mu}_K(g) = \exp\left(-\frac{1}{2}\inprod{g}{T_K g}_{L^2}\right)$

\vspace{0.5em}
$\Downarrow$
\vspace{0.5em}

\textbf{Karhunen-Lo\`eve:} $\displaystyle f \sim \mu_K \iff f = \sum_{k=1}^\infty \sqrt{\mu_k} \, \xi_k \, \phi_k$, \quad $\xi_k \iid N(0,1)$

\vspace{0.5em}
$\Downarrow$
\vspace{0.5em}

\textbf{Onsager-Machlup:} $\displaystyle \frac{\mu_K(B_\epsilon(f))}{\mu_K(B_\epsilon(0))} \xrightarrow{\epsilon \to 0} \exp\left(-\frac{1}{2}\norm{f}_{\HK}^2\right)$

\vspace{0.5em}
$\Downarrow$
\vspace{0.5em}

\textbf{Prior interpretation:} $\displaystyle p(f) \propto \exp\left(-\frac{1}{2}\norm{f}_{\HK}^2\right)$
}}
\end{center}

%%%%%%%%%%%%%%%%%%%%%%%%%%%%%%%%%%%%%%%%%%%%%%%%%%%%%%%%%%%%%%%%%%%%%%%%%%%%%%%
\section{References}
%%%%%%%%%%%%%%%%%%%%%%%%%%%%%%%%%%%%%%%%%%%%%%%%%%%%%%%%%%%%%%%%%%%%%%%%%%%%%%%

\begin{enumerate}[label={[\arabic*]}]
    \item V.I.~Bogachev. \emph{Gaussian Measures}. Mathematical Surveys and Monographs, vol.~62. American Mathematical Society, 1998.
    
    \item R.H.~Cameron and W.T.~Martin. Transformations of Wiener integrals under translations. \emph{Annals of Mathematics}, 45(2):386--396, 1944.
    
    \item O.~Zeitouni. On the Onsager-Machlup functional of diffusion processes around non-smooth curves. \emph{Annals of Probability}, 17(3):1037--1054, 1989.
    
    \item A.W.~van der Vaart and J.H.~van Zanten. Reproducing kernel Hilbert spaces of Gaussian priors. In \emph{Pushing the Limits of Contemporary Statistics: Contributions in Honor of Jayanta K. Ghosh}, IMS Collections, vol.~3, pages 200--222, 2008.
    
    \item A.~Berlinet and C.~Thomas-Agnan. \emph{Reproducing Kernel Hilbert Spaces in Probability and Statistics}. Springer, 2004.
    
    \item G.~Wahba. \emph{Spline Models for Observational Data}. CBMS-NSF Regional Conference Series in Applied Mathematics, vol.~59. SIAM, 1990.
    
    \item G.~Kallianpur. The role of reproducing kernel Hilbert spaces in the study of Gaussian processes. In \emph{Advances in Probability and Related Topics}, vol.~2, pages 49--83. Dekker, 1970.
    
    \item N.N.~Vakhania, V.I.~Tarieladze, and S.A.~Chobanyan. \emph{Probability Distributions on Banach Spaces}. Mathematics and Its Applications. D.~Reidel Publishing Company, 1987.
    
    \item L.~Gross. Abstract Wiener spaces. In \emph{Proceedings of the Fifth Berkeley Symposium on Mathematical Statistics and Probability}, vol.~2, pages 31--42, 1967.
\end{enumerate}

\end{document}
